\documentclass[12pt,a4paper]{article}

% ------------------------------------------------
% Paquetes
% ------------------------------------------------
\usepackage[utf8]{inputenc}
\usepackage[T1]{fontenc}
\usepackage[spanish,es-nodecimaldot]{babel}
\usepackage{geometry}
\usepackage{lmodern}
\usepackage{microtype}
\usepackage{enumitem}
\usepackage{array}
\usepackage{booktabs}
\usepackage{longtable}
\usepackage{xcolor}
\usepackage{hyperref}
\usepackage{fancyhdr}
\usepackage{titlesec}
\usepackage{listings}
\usepackage{tcolorbox}
\usepackage{tabularx}

\geometry{margin=2.25cm}
\hypersetup{colorlinks=true,linkcolor=blue!45!black,urlcolor=blue!55!black,
 pdftitle={Práctica semanal 01: auditoría y endurecimiento de una red de sucursal},
 pdfauthor={Laboratorio de prácticas de ingeniería de redes}}
\pagestyle{fancy}
\fancyhf{}
\lhead{Prácticas profesionales -- Redes}
\rhead{Semana 01}
\cfoot{\thepage}
\setlength{\headheight}{15pt}
\setlength{\parskip}{0.4em}
\setlength{\parindent}{0pt}

\definecolor{azul}{RGB}{25,70,115}
\definecolor{grisfondo}{RGB}{245,247,249}
\definecolor{rojo}{RGB}{145,35,35}
\titleformat{\section}{\Large\bfseries\color{azul}}{\thesection.}{0.6em}{}
\titleformat{\subsection}{\large\bfseries\color{azul}}{\thesubsection}{0.6em}{}

\lstdefinestyle{ios}{basicstyle=\ttfamily\small,backgroundcolor=\color{grisfondo},
 frame=single,rulecolor=\color{black!20},breaklines=true,columns=fullflexible,
 keepspaces=true,showstringspaces=false}
\lstset{style=ios}
\newtcolorbox{alerta}{colback=red!4,colframe=red!45!black,title=Importante,fonttitle=\bfseries}
\newtcolorbox{entrega}{colback=blue!3,colframe=azul,title=Entregables,fonttitle=\bfseries}

\begin{document}

\begin{titlepage}
\centering
\vspace*{1.3cm}
{\Large\textsc{Programa de prácticas profesionales}\par}
\vspace{0.8cm}
{\huge\bfseries Auditoría, segmentación y endurecimiento de una red de sucursal\par}
\vspace{0.5cm}
{\Large Semana 01 -- Diagnóstico y corrección de una red heredada\par}
\vspace{1.5cm}
\begin{tabular}{rl}
\textbf{Área:} & Infraestructura y comunicaciones \\
\textbf{Entorno:} & Cisco Packet Tracer \\
\textbf{Nivel:} & Intermedio inicial \\
\textbf{Tiempo estimado:} & 4--6 horas \\
\textbf{Modalidad:} & Individual \\
\end{tabular}
\vfill
\begin{tcolorbox}[width=0.9\textwidth,colback=grisfondo,colframe=black!30]
\textbf{Tu rol:} eres el ingeniero de redes en prácticas asignado a revisar una sucursal que presenta fallos intermitentes y controles de seguridad inconsistentes. No debes asumir que la configuración existente es correcta: primero recopila evidencia, identifica riesgos, plantea cambios y valida cada resultado.
\end{tcolorbox}
\vfill
{\large Documento de trabajo -- No contiene la configuración de solución\par}
\end{titlepage}

\tableofcontents
\newpage

\section{Contexto de la organización}

La empresa ficticia \textbf{NorteLogística} opera una sucursal administrativa y de almacén. La red fue ampliada por distintos técnicos durante varios meses; parte de la configuración está documentada, pero no existe un inventario confiable de cambios. La empresa no puede detener toda la operación durante el diagnóstico.

Actualmente se reportan los siguientes síntomas:
\begin{itemize}[leftmargin=1.2cm]
 \item Algunos equipos reciben direcciones IP, pero no todos pueden acceder a los servicios internos.
 \item El personal de administración puede llegar a servidores que no necesita utilizar.
 \item El equipo de soporte desconoce qué ocurrirá si falla la puerta de enlace principal.
 \item Se encontró un puerto de switch activo en una zona de visitantes.
 \item Algunos dispositivos de red todavía aceptan administración mediante Telnet.
 \item Se sospecha que un enlace entre switches transporta VLAN innecesarias.
\end{itemize}

La dirección solicita restaurar la conectividad requerida sin permitir acceso indiscriminado entre áreas. Los cambios deben ser verificables y reversibles. No se autoriza conectar esta práctica a una red real ni utilizar credenciales reales.

\subsection{Áreas y requisitos del negocio}
\begin{longtable}{p{2.6cm}p{3.4cm}p{7.0cm}}
\toprule
\textbf{Área} & \textbf{VLAN / subred} & \textbf{Necesidades y restricciones} \\
\midrule
Administración & VLAN 10 -- 192.168.10.0/24 & Acceso a DNS y a la aplicación interna HTTPS. Puede usar ping para diagnosticar interfaces de red, pero no necesita iniciar sesiones HTTP hacia los servidores. \\
Operaciones & VLAN 20 -- 192.168.20.0/24 & Acceso a DNS, DHCP y aplicación interna HTTPS. No debe acceder a la VLAN de administración. \\
Invitados & VLAN 30 -- 192.168.30.0/24 & Debe recibir direccionamiento dinámico. No debe acceder a las redes privadas internas de la empresa. No es necesario simular Internet real. \\
Servidores & VLAN 40 -- 192.168.40.0/24 & Aloja el servidor DNS y el servidor web de prueba. Solo estarán disponibles los servicios expresamente requeridos. \\
Gestión de red & VLAN 99 -- 192.168.99.0/24 & Acceso de administración restringido al equipo de soporte. No se debe permitir Telnet. \\
\bottomrule
\end{longtable}

\section{Topología y alcance}

Construye la topología en Packet Tracer con dos switches multicapa compatibles con SVI, HSRP y routing, dos switches de acceso, dos routers y dos servidores. Puedes usar switches 3560 para distribución y 2960 para acceso, o modelos equivalentes que existan en tu versión. Si un modelo no soporta una función, documenta la limitación y selecciona un modelo compatible del simulador.

\begin{center}
\begin{tcolorbox}[width=0.94\textwidth,colback=grisfondo,colframe=black!30]
\centering
\textbf{PC-ADM} (VLAN 10) \quad \textbf{PC-OPS} (VLAN 20) \quad \textbf{PC-GUEST} (VLAN 30)\\[0.25cm]
\(\downarrow\) \hspace{2.4cm} \(\downarrow\) \hspace{2.4cm} \(\downarrow\)\\[0.2cm]
\textbf{SW-ACC1} \hspace{3.0cm} \textbf{SW-ACC2}\\[0.25cm]
{\small Enlaces trunk hacia ambos switches de distribución}\\[0.2cm]
\textbf{DIST1} \(\Longleftrightarrow\) \textbf{DIST2}\\
\small EtherChannel con LACP entre los switches de distribución\\[0.25cm]
\textbf{DIST1 / DIST2:} SVIs, HSRP y routing entre VLAN\\[0.25cm]
\textbf{DIST1} \(\longrightarrow\) \textbf{R-EDGE} \(\longleftrightarrow\) \textbf{R-ISP-SIM}\\
\small Enlace de tránsito para practicar OSPF\\[0.25cm]
\textbf{Servidor DNS} y \textbf{Servidor WEB} conectados a VLAN 40.
\end{tcolorbox}
\end{center}

El diagrama es lógico y no asigna números de interfaz. Selecciona puertos, registra la asignación y mantenla consistente. No conectes enlaces paralelos entre switches sin configurar correctamente EtherChannel o sin entender cómo STP tratará los enlaces.

\subsection{Plan de direccionamiento previsto}
\begin{longtable}{p{3.1cm}p{3.0cm}p{6.9cm}}
\toprule
\textbf{Elemento} & \textbf{Dirección / red} & \textbf{Observaciones} \\
\midrule
VLAN 10 & 192.168.10.0/24 & Gateway virtual sugerido: 192.168.10.1 \\
VLAN 20 & 192.168.20.0/24 & Gateway virtual sugerido: 192.168.20.1 \\
VLAN 30 & 192.168.30.0/24 & Gateway virtual sugerido: 192.168.30.1 \\
VLAN 40 & 192.168.40.0/24 & Gateway virtual sugerido: 192.168.40.1 \\
VLAN 99 & 192.168.99.0/24 & Gateway virtual sugerido: 192.168.99.1 \\
Servidor DNS & 192.168.40.10/24 & Nombre de prueba: \texttt{dns.nortelogistica.test} \\
Servidor WEB & 192.168.40.20/24 & Nombre de prueba: \texttt{portal.nortelogistica.test} \\
Enlace de tránsito & 10.0.0.0/30 & Entre R-EDGE y R-ISP-SIM \\
\bottomrule
\end{longtable}

Define y documenta las direcciones físicas de los switches de distribución, los gateways virtuales HSRP, los enlaces de Capa 3 y las direcciones de gestión de los switches de acceso. Reserva direcciones fuera de los rangos DHCP para gateways, servidores e infraestructura. Decide rangos DHCP razonables para las VLAN de usuarios y justifica las exclusiones.

\subsection{Estado inicial de la práctica}
No se proporciona un archivo \texttt{.pkt} inicial: debes construir la topología y preparar un estado de trabajo que represente una red heredada. Para simular el relevo de otro técnico, deja inicialmente sin implementar las políticas de seguridad y los servicios que se describen en las órdenes de trabajo. Además, incorpora al menos cuatro anomalías controladas en tu estado inicial, seleccionadas de esta lista:
\begin{itemize}[leftmargin=1.2cm]
 \item Un puerto de usuario asignado a una VLAN equivocada.
 \item Un trunk que no permite una de las VLAN requeridas.
 \item Un ámbito DHCP con gateway o DNS incorrecto.
 \item Una ACL que bloquea HTTPS requerido o permite tráfico innecesario.
 \item Una línea VTY que todavía acepta Telnet.
 \item Un puerto de usuario sin Port Security.
 \item Un EtherChannel con parámetros incompatibles en un extremo.
 \item Una configuración HSRP sin el comportamiento de recuperación que necesita el negocio.
\end{itemize}
Anota en privado qué anomalías sembraste para poder confirmar después si las detectaste todas. No las corrijas al construir el estado inicial: primero completa el inventario y las pruebas de diagnóstico. El informe debe distinguir entre un fallo que observaste y uno que solo suponías.

\section{Fase 1 -- Inspección antes de modificar}

Antes de configurar, realiza un inventario de dispositivos y enlaces. En cada equipo, utiliza los comandos de inspección aplicables. Algunos ejemplos:
\begin{lstlisting}
show running-config
show startup-config
show ip interface brief
show interfaces status
show vlan brief
show interfaces trunk
show etherchannel summary
show spanning-tree
show ip route
show ip ospf neighbor
show standby brief
show access-lists
show port-security
show logging
\end{lstlisting}

No todos los comandos aplican a todos los dispositivos. Registra cuáles están disponibles y cuáles no.
\begin{enumerate}[leftmargin=1.2cm]
 \item Identifica interfaces conectadas, su estado y el tipo de enlace.
 \item Comprueba qué VLAN existen y cuáles están permitidas por cada trunk.
 \item Determina si EtherChannel está agregado y qué protocolo utiliza.
 \item Identifica las puertas de enlace configuradas en los clientes y el estado de HSRP.
 \item Comprueba las rutas instaladas y el estado de OSPF, si existe.
 \item Revisa las ACL en orden y determina qué tráfico permiten o bloquean.
 \item Revisa si los puertos de usuario tienen Port Security y si la administración remota usa SSH.
\end{enumerate}

\textbf{Resultado esperado:} una tabla de hallazgos con al menos cinco observaciones, respaldadas por salidas de comandos, pruebas o capturas. Guarda tu diagnóstico inicial antes de modificar la configuración.

\section{Fase 2 -- Diseñar antes de configurar}

Prepara los siguientes artefactos:
\begin{enumerate}[leftmargin=1.2cm]
 \item Una tabla puerto--dispositivo--VLAN--modo (access, trunk o routed).
 \item Una tabla de direccionamiento con IP, máscara, gateway, DNS y asignación estática o dinámica.
 \item Una matriz de políticas con origen, destino, protocolo/puerto, acción y justificación.
 \item Una lista de cambios ordenados por riesgo, con un procedimiento de reversión para cada cambio importante.
\end{enumerate}

No basta con escribir «configurar ACL». Describe qué comunicaciones necesita cada área y cuáles deben estar prohibidas.

\section{Fase 3 -- Switching y segmentación}
\begin{enumerate}[leftmargin=1.2cm]
 \item Crea las VLAN 10, 20, 30, 40 y 99 donde corresponda.
 \item Asigna los puertos de usuario a la VLAN adecuada en modo access.
 \item Configura los trunks entre acceso y distribución, limitando las VLAN permitidas a las necesarias.
 \item Configura un EtherChannel entre DIST1 y DIST2 mediante LACP. Verifica que los parámetros de ambos extremos sean compatibles.
 \item Configura PortFast y BPDU Guard únicamente en puertos de borde que conecten dispositivos finales.
 \item Apaga los puertos sin uso y asígnalos a una VLAN aislada si el modelo lo permite.
\end{enumerate}

\section{Fase 4 -- Routing, DHCP y DNS}
\begin{enumerate}[leftmargin=1.2cm]
 \item Configura las interfaces VLAN (SVI) y el routing entre VLAN en los switches multicapa.
 \item Configura HSRP en DIST1 y DIST2 para las VLAN que requieran gateway redundante. Define prioridades y decide si usarás \texttt{preempt}; justifica el comportamiento de recuperación.
 \item Configura el enlace de Capa 3 entre R-EDGE y R-ISP-SIM, asigna las direcciones del segmento de tránsito y establece OSPF para intercambiar las rutas previstas.
 \item Decide dónde se alojará DHCP. Si el servidor DHCP está en otra subred, configura el relay que soporte el equipo y explica en qué interfaz debe operar.
 \item Configura ámbitos DHCP para las VLAN de usuarios. Excluye direcciones reservadas y entrega gateway y DNS correctos.
 \item Configura DNS con un registro que resuelva \texttt{portal.nortelogistica.test} a la IP del servidor web.
 \item Habilita el servicio HTTP/HTTPS de prueba disponible en Packet Tracer y comprueba el acceso mediante nombre DNS, no solo por IP.
\end{enumerate}

R-ISP-SIM representa una red externa simulada; no se necesita conectividad real a Internet.

\section{Fase 5 -- Seguridad}
\begin{enumerate}[leftmargin=1.2cm]
 \item Configura Port Security en puertos de usuario. Define máximo de MAC, método de aprendizaje y modo de violación. Registra qué sucede al conectar una MAC no autorizada.
 \item Configura acceso remoto por SSH y deshabilita Telnet en las líneas VTY de los equipos que lo soporten.
 \item Restringe la administración a la VLAN 99 o a los equipos autorizados dentro de ella. Evita bloquear tu propio acceso.
 \item Diseña e implementa ACL extendidas para que Administración y Operaciones solo usen los servicios internos necesarios.
 \item Impide que Invitados inicie comunicaciones hacia las redes privadas internas.
 \item Conserva únicamente las excepciones necesarias para DHCP, DNS, HTTPS y las pruebas de diagnóstico autorizadas.
 \item Si Packet Tracer lo permite, investiga DHCP Snooping y Dynamic ARP Inspection. Documenta los puertos confiables/no confiables y prueba el comportamiento.
\end{enumerate}

\begin{alerta}
Las ACL IPv4 se evalúan en orden, de arriba hacia abajo, y terminan en una denegación implícita. Una regla general de permiso antes de una denegación específica puede volver ineficaz la política. Verifica el orden, la interfaz y la dirección (\texttt{in}/\texttt{out}) donde aplicas cada ACL.
\end{alerta}

\section{Plan de pruebas obligatorio}

No declares terminada la práctica únicamente porque los comandos se hayan aceptado. Registra el resultado esperado, el observado y la evidencia de cada prueba.
\begin{longtable}{p{0.6cm}p{5.1cm}p{4.2cm}p{3.0cm}}
\toprule
\textbf{\#} & \textbf{Prueba} & \textbf{Resultado esperado} & \textbf{Evidencia} \\
\midrule
1 & Consultar VLAN y trunks en cada switch. & VLAN correctas y trunks con la lista prevista. & Comandos \texttt{show}. \\
2 & Revisar EtherChannel en ambos extremos. & Canal operativo y puertos miembros agregados. & \texttt{show etherchannel summary}. \\
3 & Obtener IP por DHCP en PC-ADM, PC-OPS y PC-GUEST. & IP, máscara, gateway y DNS del ámbito correcto. & \texttt{ipconfig}. \\
4 & Hacer ping al gateway virtual desde cada VLAN. & Respuesta cuando la política lo permite. & Captura de resultado. \\
5 & Resolver \texttt{portal.nortelogistica.test}. & Devuelve la IP del servidor web. & \texttt{nslookup} si está disponible o prueba DNS en PT. \\
6 & Abrir el portal por nombre DNS desde Administración y Operaciones. & HTTPS permitido según política. & Navegador y prueba de conectividad. \\
7 & Intentar HTTP hacia el portal desde VLAN internas. & Bloqueado si no está autorizado expresamente. & Navegador, ACL y contadores. \\
8 & Desde Invitados, intentar acceder a DNS y servidor web internos. & Solo servicios explícitamente autorizados deben funcionar. & Pruebas DNS, HTTP/HTTPS y ping. \\
9 & Probar conectividad entre Administración y Operaciones. & No hay acceso indiscriminado entre VLAN. & Pings y pruebas de aplicación. \\
10 & Verificar OSPF y rutas aprendidas. & Adyacencia y rutas esperadas. & \texttt{show ip ospf neighbor}, \texttt{show ip route}. \\
11 & Apagar temporalmente la interfaz o equipo HSRP activo. & El otro equipo asume el rol activo y se recupera conectividad. & \texttt{show standby brief}. \\
12 & Restaurar el equipo principal. & El rol activo coincide con la política de prioridad/preemption. & Estado HSRP antes y después. \\
13 & Conectar una MAC distinta a un puerto protegido. & Se ejecuta la acción de violación configurada. & \texttt{show port-security interface}. \\
14 & Probar SSH desde un host autorizado y Telnet. & SSH funciona; Telnet se rechaza. & Cliente SSH/Telnet y VTY. \\
15 & Intentar administrar desde un host no autorizado. & Se rechaza administración sin afectar servicios legítimos. & Prueba y ACL. \\
16 & Revisar registros y contadores de ACL. & Pruebas reflejadas en contadores o registros disponibles. & \texttt{show access-lists}, \texttt{show logging}. \\
\bottomrule
\end{longtable}

Algunas versiones de Packet Tracer no implementan todas las funciones de equipos reales ni muestran los mismos contadores. Si una prueba no es viable, anota la limitación concreta y cómo la comprobarías en hardware real. No declares exitosa una prueba que no pudiste ejecutar.

\section{Criterios de aceptación}
La práctica se considera terminada cuando:
\begin{itemize}[leftmargin=1.2cm]
 \item Las VLAN y trunks coinciden con el diseño y no transportan VLAN innecesarias.
 \item EtherChannel funciona como enlace lógico y no provoca bucles.
 \item Los clientes reciben direccionamiento correcto y resuelven el nombre DNS del portal.
 \item OSPF aprende las rutas previstas.
 \item HSRP proporciona redundancia y la recuperación coincide con la política definida.
 \item Las ACL permiten flujos necesarios y bloquean los no autorizados.
 \item Port Security y SSH funcionan según la configuración.
 \item Se ejecutaron las pruebas y se registraron éxitos y fallos.
 \item La configuración final está guardada y el informe documenta las limitaciones del simulador.
\end{itemize}

\section{Entregables}
\begin{entrega}
Entrega un informe en PDF y el archivo de Packet Tracer (\texttt{.pkt}):
\begin{enumerate}[leftmargin=1.2cm]
 \item \textbf{Diagnóstico inicial:} topología, inventario, estado de enlaces y al menos cinco hallazgos con evidencia.
 \item \textbf{Diseño:} diagrama final, tabla de VLAN, direccionamiento, puertos y política de tráfico.
 \item \textbf{Cambios realizados:} resumen de configuración por dispositivo y explicación de decisiones relevantes.
 \item \textbf{Pruebas:} tabla con los 16 casos, resultado observado, esperado y evidencia.
 \item \textbf{Seguridad:} amenazas mitigadas, riesgos restantes y funciones que no pudiste implementar en Packet Tracer.
 \item \textbf{Reversión:} cómo desharías al menos tres cambios críticos si afectaran la operación.
 \item \textbf{Reflexión técnica:} tres errores o dificultades, cómo los diagnosticaste y qué aprendiste.
\end{enumerate}
\end{entrega}

\section{Registro de hallazgos}
Utiliza una tabla como esta durante la auditoría:
\begin{longtable}{p{2.4cm}p{2.8cm}p{2.7cm}p{3.1cm}p{1.5cm}}
\toprule
\textbf{Dispositivo} & \textbf{Hallazgo} & \textbf{Evidencia} & \textbf{Acción propuesta} & \textbf{Prioridad} \\
\midrule
Ejemplo: SW-ACC1 & VLAN incorrecta en un puerto de usuario & \texttt{show vlan brief} & Corregir asignación tras validar el puerto & Alta \\
\bottomrule
\end{longtable}

Prioriza hallazgos por impacto en confidencialidad, integridad, disponibilidad y alcance del fallo. No corrijas un síntoma sin entender su causa.

\section{Reglas de trabajo}
\begin{enumerate}[leftmargin=1.2cm]
 \item No borres toda la configuración para empezar de cero: diagnostica y determina qué puede reutilizarse.
 \item Guarda una copia de las configuraciones antes de cambios importantes.
 \item Aplica cambios pequeños y verifica después de cada fase.
 \item Si pierdes conectividad, recupera el acceso antes de continuar.
 \item No confundas «el comando fue aceptado» con «el requisito quedó satisfecho».
 \item Si una función no está soportada, registra evidencia y no inventes resultados.
 \item Utiliza direcciones y credenciales ficticias exclusivamente para el laboratorio.
\end{enumerate}

\vfill
\begin{center}
\textcolor{azul}{\textbf{Fin de la orden de trabajo -- Semana 01}}\\
\small El objetivo no es introducir la mayor cantidad de comandos, sino producir una red verificable, segmentada y mantenible.
\end{center}
\end{document}
